\documentclass[11pt]{article}
\usepackage[margin=1in]{geometry}
\usepackage{amsmath,amssymb,amsthm}
\usepackage{enumitem}

\newtheorem{theorem}{Theorem}
\newtheorem{proposition}[theorem]{Proposition}
\newtheorem{corollary}[theorem]{Corollary}
\newtheorem{remark}[theorem]{Remark}

\title{d5\_236: A finite-defect splice principle for the remaining graph-side theorem \(G1\)}
\author{}
\date{2026-03-21}

\begin{document}
\maketitle

\section{Purpose}
After the graph-side split, the only live object is \(G1\): splice the already
closed color-$4$ branch into a colorwise package while preserving
\begin{enumerate}[label=\textup{(\arabic*)},leftmargin=2.8em]
\item pointwise outgoing exhaustiveness, and
\item colorwise incoming uniqueness.
\end{enumerate}
The present note packages the exact finite-defect theorem kernel needed for that
step.
It is deliberately coordinate-free.

\section{The splice principle}
Let \(X\) be a finite set.
For each color \(c\in\{0,1,2,3,4\}\), let
\[
 f_c\colon X\to X
\]
be a permutation.
Think of \((f_0,\dots,f_4)\) as the already valid baseline colorwise package.
Let \(D\subseteq X\) be a defect set, and for each color let
\[
 g_c\colon X\to X
\]
be a modified map such that
\[
 g_c(x)=f_c(x)\qquad (x\notin D).
\]

\begin{proposition}[one-color finite-defect splice]
\label{prop:one-color-splice}
Fix one color \(c\).
Assume:
\begin{enumerate}[label=\textup{(\arabic*)},leftmargin=2.8em]
\item \(g_c\vert_D\) is injective;
\item the defect image is preserved setwise:
      \[
      g_c(D)=f_c(D).
      \]
\end{enumerate}
Then \(g_c\) is a permutation of \(X\).
\end{proposition}

\begin{proof}
Because \(f_c\) is a permutation, its image on \(X\setminus D\) is exactly
\[
 f_c(X\setminus D)=X\setminus f_c(D).
\]
Since \(g_c=f_c\) on \(X\setminus D\), we also have
\[
 g_c(X\setminus D)=X\setminus f_c(D)=X\setminus g_c(D).
\]
Thus the images of \(D\) and \(X\setminus D\) under \(g_c\) are disjoint and
cover all of \(X\).
By injectivity on \(D\) and because \(g_c\) agrees with the permutation \(f_c\)
on \(X\setminus D\), the map \(g_c\) is injective on each part and has disjoint
images on the two parts.
Hence \(g_c\) is bijective.
\end{proof}

\begin{corollary}[colorwise splice criterion]
\label{cor:colorwise-splice}
Suppose for every color \(c\) the hypotheses of
Proposition~\ref{prop:one-color-splice} hold.
Then every modified color map \(g_c\) is a permutation.
So to preserve colorwise incoming uniqueness, one only needs a finite defect
certificate on the images of the defect sets.
\end{corollary}

\section{Outgoing exhaustiveness is also a defect-local check}
Assume in addition that the baseline package satisfies pointwise outgoing
exhaustiveness:
for every \(x\in X\), the five chosen outgoing directions form the full set
\(\{0,1,2,3,4\}\).
If the modified package agrees with the baseline outside a defect set \(D\),
then pointwise outgoing exhaustiveness can fail only on \(D\).
Therefore the two graph-side obligations split cleanly:
\begin{enumerate}[label=\textup{(\arabic*)},leftmargin=2.8em]
\item \textbf{Permutation side:} verify the image-preserving splice condition of
      Corollary~\ref{cor:colorwise-splice};
\item \textbf{Selector side:} verify pointwise outgoing exhaustiveness on the
      finite defect set \(D\).
\end{enumerate}

\begin{theorem}[finite-defect form of \(G1\)]
\label{thm:G1-finite-defect}
Let \((f_0,\dots,f_4)\) be a valid baseline colorwise package on a finite set
\(X\), and let \((g_0,\dots,g_4)\) be a modified package that agrees with the
baseline outside a defect set \(D\subseteq X\).
Assume:
\begin{enumerate}[label=\textup{(\arabic*)},leftmargin=2.8em]
\item for each color \(c\), the modified defect map \(g_c\vert_D\) is injective;
\item for each color \(c\), the defect image is preserved setwise:
      \[
      g_c(D)=f_c(D);
      \]
\item pointwise outgoing exhaustiveness holds for every \(x\in D\).
\end{enumerate}
Then \((g_0,\dots,g_4)\) is again a valid colorwise package:
all five maps are permutations and outgoing exhaustiveness holds pointwise on
all of \(X\).
\end{theorem}

\begin{proof}
By Corollary~\ref{cor:colorwise-splice}, each \(g_c\) is a permutation.
Outgoing exhaustiveness is unchanged on \(X\setminus D\) because the modified
package agrees there with the baseline package.
It holds on \(D\) by hypothesis.
\end{proof}

\section{Interpretation for the current manuscript}
This theorem does \emph{not} yet close \(G1\) by itself, because the manuscript
still needs the actual defect data relating the corrected selector baseline to
the closed color-$4$ branch.
But it identifies the exact theorem kernel that remains:
\begin{itemize}[leftmargin=2.4em]
\item surface the baseline package and the closed color-$4$ branch on a common
      state space;
\item identify the finite defect set where they differ;
\item verify image-preserving splice for each color on that defect set;
\item verify pointwise outgoing exhaustiveness on that same defect set.
\end{itemize}
So after the closure of \(G2\), the residual content of \(G1\) is a concrete
finite-defect splice certificate.

\end{document}
